% ===========================================================================
\section{Completeness: What Is Estimated and What Is Not}
\label{sec:completeness}
% ===========================================================================

IPCC requires that an inventory account for every source category it might contain, using
notation keys where a number is absent: NE for not estimated, NO for not occurring, IE for
included elsewhere, NA for not applicable. The requirement is easy to satisfy formally and
rarely satisfied usefully, because a table of NE entries with no magnitudes tells a reader
that something is missing without telling them whether it matters.

We report notation keys over \num{138} category--gas combinations spanning \num{72} IPCC codes,
and where a source exists that can bound a gap, the completeness table carries the
magnitude. Table~\ref{tab:completeness} summarises; the full table is generated by the engine
and versioned with it.

\begin{table}[htbp]
% the counts are value macros from four stored queries (get_completeness per key);
% see paper_national/queries/paper_artefacts.csv, tab_completeness
% GENERATED by scripts/make_paper_artefacts.py from paper_artefacts.csv: query=get_completeness params_sha256=e50a8c6f1dc58283c44410b437f2dabccbca35a8c1469b450f95027c38395b54 result_sha256=ccea93bc6f19cc4d292b65efd84418ffbc234603a35235dab33829b465812317 release=PakEmis 1.0.0rc68 receipt=1c13d46a726be57378d2f8d422986afe;bf31bd881e966273008b12df56d94417;dd745ba7decdb14d1ba678da8ec03181;5c18389f35dbe70e0db988b024ed2fc4;2c102cfc2a2c35cbd556dedd4cbc2536 -- do not edit
\providecommand{\genCompletenessCountE}{}\renewcommand{\genCompletenessCountE}{54}% receipt 1c13d46a726be57378d2f8d422986afe
\providecommand{\genCompletenessCountNE}{}\renewcommand{\genCompletenessCountNE}{71}% receipt bf31bd881e966273008b12df56d94417
\providecommand{\genCompletenessCountIE}{}\renewcommand{\genCompletenessCountIE}{8}% receipt dd745ba7decdb14d1ba678da8ec03181
\providecommand{\genCompletenessCountNO}{}\renewcommand{\genCompletenessCountNO}{4}% receipt 5c18389f35dbe70e0db988b024ed2fc4
\providecommand{\genCompletenessCountNA}{}\renewcommand{\genCompletenessCountNA}{1}% receipt 2c102cfc2a2c35cbd556dedd4cbc2536

\centering
\small
\begin{tabular}{@{}clp{0.52\linewidth}@{}}
\toprule
\textbf{Key} & \textbf{Count} & \textbf{Meaning in this inventory} \\
\midrule
E  & \genCompletenessCountE{} & Estimated and reported \\
NE & \genCompletenessCountNE{} & Not estimated. Each carries a justification, and those a comparator can bound
          carry a magnitude \\
IE & \genCompletenessCountIE{} & Included elsewhere: the cross-sector transfers of Section~\ref{sec:closures},
          household electricity counted at the power station, biogenic \COtwo{} held as a
          memo, and the transport boundary's declared inclusions---stationary
          agricultural plant, captive diesel and construction machinery reported inside
          the road and off-road controls at those categories' factors. (Railway
          non-traction fuel left this key: it is an estimated row at its own
          1A2/1A4c-unresolved category rather than a declaration inside 1A3c) \\
NO & \genCompletenessCountNO{} & Not occurring: managed savanna burning, other agricultural sources, direct
          combustion on urban electric rail, and electric road vehicles \\
NA & \genCompletenessCountNA{} & Not applicable: soda ash production (2.B.7), where the category exists in
          the reporting hierarchy and the process does not arise in Pakistan---distinct
          from NO, which records an activity that does not occur, and from NE, which
          records one that does and is unquantified \\
\bottomrule
\end{tabular}
\caption{Completeness across \num{138} category--gas combinations and \num{72} IPCC codes.
The generated table carries a justification for every non-E entry.}
\label{tab:completeness}
\end{table}

\subsection{The declared holes, in order of size}

\textbf{Manufacturing industries and construction (1A2), roughly \SI{15.85}{\mega\tonne}
remaining --- halved in this release by the solid-fuel activation.} EDGAR puts Pakistani
industrial fuel combustion near \SI{44}{\mega\tonne}. This inventory reports
\SI{5.55}{\mega\tonne} of captive self-generation
(screened by configuration and fuel class, with a bracket and a status per row, in
\texttt{outputs/industry\_\allowbreak captive\_\allowbreak screening\_\allowbreak 2024.csv})
and \SI{5.61}{\mega\tonne} of industrial
process gas, both from OGRA's sector-wise gas table, \SI{3.40}{\mega\tonne} of industrial
liquid fuels from OCAC's Total Industry customer class, and---new in this release, through the
category's written activation gate---\SI{2.78}{\mega\tonne} of brick-kiln and
\SI{10.80}{\mega\tonne} of cement-sector coal from the HDIP coal-by-consumer series
republished in the Finance Division's Statistical Supplement. About \SI{15.85}{\mega\tonne}
is outstanding against \SI{29}{} in the previous edition, \SI{33}{} before that and
\SI{38}{} before that.

Both of the recent additions closed an error rather than a data gap, and they closed the same
error twice. This inventory first recorded that Pakistan publishes no sector-disaggregated
industrial fuel consumption. It does: OGRA reports gas by sector, in the table that was already
the source of this inventory's residential control total, its LPG total, its kerosene total and
the captive series. The register was then rewritten to say the remainder was specifically
non-gas and required a paid source. That was also wrong: OCAC's Total Industry customer
class---furnace oil, diesel, LDO, kerosene and gasoline---had been in the repository since the
transport sector first ingested the same annual panel for its road control total, and no module
read those rows. What was missing was not the statistic but the reading of it, twice, in the
same category---which is the general finding of Section~\ref{sec:discussion}'s account of the
gap register, and the reason the register now distinguishes a datum that was never acquired
from one that was acquired and never extracted.

The activation is the graduation the previous edition named. Its screenings had read the old
\SI{29}{\mega\tonne} as largely interpreted---kiln fuels, mostly---and said what would
graduate them: the HDIP industrial coal breakdown. That series arrived (the Finance Division's
Statistical Supplement republishes it, Table 14.1, with categories that exhaust supply), and
the two large claimants passed through the category's written gate as reference streams priced
on each year's own coal. Two findings came with them. The measured cement-sector coal
(\SI{10.80}{\mega\tonne} on \num{4391}~kt) sits \emph{below} the superseded transfer
model's \SIrange{11.99}{15.65}{\mega\tonne} bracket---either the transferred heat intensity
overstates Pakistani kilns or part of the kiln heat is not coal; the model is demoted to a
comparator with that finding attached, not reconciled away. And the brick column is a
residual-style bucket in HDIP's own accounting---supply exhausted means whatever industrial
coal is not called power, cement or household lands in it---which is priced as the stream's
\SI{30}{\percent} activity band rather than footnoted.

The codes in this table are IPCC 2006 throughout, from Volume~1 Chapter~8 Table~8.2. An
earlier release coded sector~3 to the UNFCCC common reporting tables and sector~4 to the
2006 Guidelines under one column, which made a join on the code string unsafe; the scheme
and its crosswalk to the reporting tables, including the two categories that have no
counterpart in either direction, are published as
\texttt{outputs/pakemis\_category\_crosswalk.csv}.

That characterisation is now measured rather than asserted, and the measurement says the
band is too narrow. \texttt{outputs/pakemis\_coal\_allocation\_identity.csv} puts the
national coal supply beside the consumer columns summed: over the \num{14} years both tables
cover they close to within \SI{1}{\kilo\tonne} in \num{13}, far closer than two independent
statistics would agree, which is the signature of a constructed identity rather than a
balance. Inside it the split is a redistribution, and over FY2019--20 to FY2023--24 the
constraint binds: the total moves \SI{438}{\kilo\tonne} while the power column gains
\SI{8299}{\kilo\tonne} and brick kilns plus cement lose \SI{8739}{\kilo\tonne}. Every
tonne the Thar ramp added to the power column came out of another column, so the brick
series' \SI{86.2}{\percent} fall over those years is not independent evidence about brick
activity. Over the whole held span the same total grew by \SI{17001}{\kilo\tonne}, so the
redistribution reading belongs to the recent window, and the product publishes both.

The consequence for this inventory is that the priced band does not reach the only
independent bound the repository holds. At \SI{1129}{\kilo\tonne} the booked column spans
\SIrange{790}{1468}{\kilo\tonne} under its \SI{30}{\percent} band, while the brick
fleet---counted once in late 2020 at \num{10394} kilns on one reading and \num{9486} on the
other---bounds its coal at \SI{3301}{\kilo\tonne} and \SI{3095}{\kilo\tonne} even on the
shortest firing season either source declares. The band's upper end is a factor of
\num{2.25} and \num{2.11} short of those. Read in kilns at the survey's own
\SI{328.2}{\tonne} per kiln-year, the booked column supports \num{3440} firing kiln-years
at the reference year against \num{24936} in FY2019--20---from more than twice the counted
fleet to about a third of it, for a fleet counted once. \texttt{outputs/pakemis\_coal\_sector\_reallocation.csv} gives that shortfall a candidate
source, from a direction with no connection to kiln counting. Because HDIP's consumption is
its supply, replacing the power column with the stations' own metered coal changes the
industrial residual by the same tonnage the other way---at the reference year by
\SI{3321}{\kilo\tonne}, worth \SI{8.14}{\mega\tonne}~\COe{} against the power lane's
\SI{-18.03}{\mega\tonne}. Booked, the brick column sits below the fleet's own floor in
\num{3} of the \num{9} years; with that residual it is inside the bound in all \num{9}. The
claim this supports is modest, because six years were already inside and the bound is wider
than the quantity in dispute: the reallocation never worsens the brick position and repairs
it where it was broken. Which non-power consumer took the coal is not resolved.

The fleet estimate is still not booked
and should not be: it has no time axis, which is precisely why it was published beside the
column rather than in place of it. What changes is that the brick stream's uncertainty is now
known to be understated, with the factor stated.

What remains outstanding is narrower
again: petroleum coke (never in the coal table; still the Comtrade trade-screen lane), captive
coal (a \SIrange{0.28}{1.31}{\mega\tonne} configuration, deliberately \emph{not} activated
because it sits inside the brick residual and would double count), and process bagasse, still
bounded rather than estimated: the entire cane crop's bagasse reaches only
\SI{0.39}{\mega\tonne} of countable CH$_4$ and N$_2$O
(\emph{bagasse\_ceiling\_in\_scope}). Given
that every narrowing so far came from reading the repository rather than from acquiring
anything, the remaining \SI{15.85}{\mega\tonne} is better read as an upper bound on what is absent than
as a measurement of it.

One row of that same table is deliberately \emph{not} taken, and the reason is a method,
not a materiality judgement. Fertiliser draws \SI{756}{} MMCFD, more than twice general
industry and second only to power. The IPCC ammonia method allocates a plant's \emph{total}
gas requirement---feedstock and process fuel together---to the ammonia category, 2B1, and
removes that gas from Energy so the same molecules are never counted twice; only separately
metered auxiliary combustion outside the ammonia process boundary stays in 1A2. NFDC's
monthly plant annex puts the split at \SI{83.5}{\percent} feedstock and
\SI{16.5}{\percent} process fuel in 2023--24, the reference year itself, so a sixth of the
fertiliser row is burned---and the burned sixth belongs to 2B1 anyway, because the method
draws the boundary around the process rather than around the flame. The same annex puts
the sector at \SI{751.9}{} MMCFD against OGRA's \SI{756}{}, two independent Pakistani
records of one fiscal year \SI{0.5}{\percent} apart: an industry body's plant meters
against a regulator's utility returns. An earlier edition of this section said almost none
of the row was combustion, which was wrong twice over: a sixth of it is, and that was never
why the row is excluded. The exclusion stands on the method: booking the row as industrial
fuel would inflate the category three and a half fold and double-count against the ammonia
accounting.

That the estimated part is \emph{measured} rather than modelled is worth noting against the
usual pattern for a gap this size: OGRA publishes the gas annually, OCAC publishes the liquid
products monthly and annually, and a 1,105-plant register assembled from the two gas utilities'
captive fleets and the NEPRA licence list reconciles against the national figure to within
\SI{0.01}{} MMCFD. The dominant uncertainty on the gas half is not the activity but the heating
value, which nothing read here pins: \SI{900}{} against \SI{1000}{} Btu per standard cubic foot
moves the gas rows from \SI{5.05}{\mega\tonne} to \SI{5.61}{\mega\tonne} and the industrial
category as a whole from \SI{22.02}{\mega\tonne} to \SI{22.58}{\mega\tonne}, and both are
reported. The liquid rows do not move with it: they are metered in tonnes and converted by
calorific value, so the difference between the two figures is a direct measure of how much of
this category still rests on an unpinned gas assumption.

\textbf{Fluorinated gases (2F), \SI{11.54}{\mega\tonne} on EDGAR's estimate.} Refrigeration
and air-conditioning substitutes for ozone-depleting substances. If EDGAR is approximately
right, Pakistani F-gases are larger than the waste sector's solid-waste disposal category and
a third the size of waste overall. They are absent from the totals---but the absence is a
\emph{machine-declared state}, not a sentence: an HFC equipment-bank model is implemented
(\texttt{pkinv.sectors.ippu.fgases}) and gated off because only an aggregate consumption
control is held, with no bank-resolving import series. The same is true across the
thirty-five non-cement IPPU category rows: every one carries an NE state in the register,
and the chemical categories additionally publish
\texttt{outputs/ippu\_chemical\_activation\_status\_2024.csv} (which estimator exists,
what activity it holds, and the named evidence that would activate it---ammonia's
carbon-ledger estimator waits on complete plant feedstock and CO$_2$-recovery data, nitric
acid's abatement estimator on national acid output) and
\texttt{outputs/ippu\_chemical\_tier\_gate\_2024.csv} (why each tier is not ready). The
completeness ask is answered by those shipped, tested states; what none of them can supply
is the species- and application-resolved activity behind them. The customs half is now
processed (rc30): UN~Comtrade's HS-2022 HFC lines give net substance imports through
customs as a fourth typed screening lane, corroborating the UNEP supply aggregate without
resolving the mixture split---which is exactly what the National Ozone Unit request
remains for.

\textbf{Captive and industrial self-generation (1A2 partial).} IPCC places
self-generation in 1A2 on Table~2.3 factors. NEPRA's regulatory scope excludes it, which is
most of the reason Ember's Pakistani electricity total exceeds NEPRA's. It is inside the
\SI{44}{\mega\tonne} 1A2 hole rather than additional to it.

\textbf{Fugitive activity detail (1B2, partial).} Transmission, distribution, processing,
refining and post-meter sources are now included under the 2019 Refinement. The remaining
gaps are exploration, abandoned wells, episodic releases and the facility characteristics
needed to replace higher/default practice factors: pipeline material and pressure, LDAR,
compressor seals, plant throughput and well status.

\textbf{Industrial wastewater (4D2).} Needs industrial output by sector and organic loading
per unit of output, neither available nationally. Domestic wastewater is reported;
industrial is not, and given that domestic alone is \SI{5.5}{\percent} of the inventory, this
gap is not small.

\textbf{Water-borne navigation (1A3d), domestic and international.} The domestic activity
figure now exists in this repository---OCAC's sectoral table carries a Domestic Bunkers row
(\SI{6.9}{\kilo\tonne} of diesel in FY2021--22, read by the transport fuel ledger)---so
what remains is a small reconciliation, not an absent statistic. The gap is still worth
describing precisely because it contains a trap. The national petroleum statistics carry an ``Overseas/Export'' furnace-oil line of
\SI{1.01}{\mega\tonne} that reads like bunker fuel and is not: it is exported cargo, matched
independently by commercial trade data showing \SI{1.11}{\mega\tonne} of furnace-oil exports
in the same period. Karachi's bunkering operation began only in 2026. Booking that line as
bunkers would create a memo item of roughly \SI{3}{\mega\tonne} \COtwo{} out of a commodity
export.

\textbf{Iron and steel (2C1), lime (2A2), ammonia (2B1), soda ash (2B7).} Four IPPU
categories where the emission factor is uncontroversial and the activity data are the
obstacle. The repository holds production figures near two of them---\SI{9100}{\kilo\tonne}
of \emph{total fertiliser product}, a sum over urea, DAP, CAN and the rest from the Punjab
Pocket Book, and \SI{785}{\kilo\tonne} of soda ash from the PBS quantum index, which
supersedes the \SI{651}{\kilo\tonne} an earlier edition of this section still carried---and
it is worth being precise about why neither is the process activity the method needs.

\textbf{The held figures are not the process activity the method needs.} No fertiliser
tonnage is ammonia: the IPCC method needs ammonia output itself, the feedstock carbon
entering the process, and the \COtwo{} recovered downstream for urea synthesis, and a
total-product tonnage supplies none of the three. Nitric acid
needs acid output, the abatement technology installed and its operating factor.
\emph{No magnitude is asserted for these categories here.} A previous edition of this section
said the four were together ``of order \SI{1}{\mega\tonne}'', which had no category-by-category
derivation behind it and is withdrawn; Pakistan's own filing puts the chemical industries near
\SI{19.4}{\mega\tonne} for 2022, which is the right order of magnitude for the question and
is not a substitute for an independent estimate either. They stay NE because the available
production quantities do not match the physical process boundaries the IPCC methods require,
which is a statement about evidence rather than about size. For ammonia the evidence has
now been assembled far enough to test the official figure: the NFDC review supplies
production by product and manufacturer and the plant-wise gas split, and the generated
reconciliation table compares four independent activity pathways---official filing, plant
census, product stoichiometry and gas balance. They do not agree. The filing's
\SI{11467}{\kilo\tonne} of ammonia would need roughly half again more gas than the entire
fertiliser sector draws, while embodied ammonia in the products made sits near
\SI{4500}{\kilo\tonne}; the official activity reads like a total-product or capacity
figure, and until the pathways reconcile the category stays NE.

\subsection{Notation keys as a first-class output}

Every NE entry above carries a justification in the generated completeness table, and where a
comparator can bound the gap, its magnitude. That is a deliberate inversion of the usual
practice, in which the estimated categories are documented in detail and the unestimated ones
receive a two-word note.

The reason is that a reader's most important question about an inventory reporting
\SI{445}{\mega\tonne} is not how that total was computed but what is not in it.
Section~\ref{sec:benchmarks} gives the aggregate answer---this inventory is \SI{87}{\percent}
of EDGAR's Pakistani sectors summed---and the completeness table is what makes that
\SI{13}{\percent} attributable to specific named categories rather than being a diffuse
disagreement.

Four IE entries deserve their own note because they are the only cases where a real emission
is deliberately reported at zero in the place a reader might look for it. Household
electricity---\SI{39728}{\giga\watt\hour}, nearly half of national consumption---is not in
the residential sector, because those emissions occur at the power station and are counted in
1A1a; adding them would double-count roughly a quarter of the inventory. Biogenic \COtwo{}
from dung, firewood and landfilled organic matter is a memo item outside the total, per
IPCC's biomass convention and Pakistan's own stated practice. And the two cross-sector
transfers of Section~\ref{sec:closures} are IE by construction. Each is asserted by a test:
the residential sector fails its build if an electricity stream is ever added.
